% TOP_PROBLEM: {"problemId": "problem.slice-ribbon-conjecture", "canonicalTitle": "Slice-Ribbon Conjecture", "releaseRank": 301, "primaryDomain": "geometry_topology", "primaryDomainLabel": "Geometry & topology", "displayStatus": "open_with_solved_subcases", "releaseStatus": "open_with_solved_subcases"}
% !TeX program = lualatex
% ATTEMPT_STATUS: unresolved
\documentclass[11pt]{article}
\usepackage[margin=25mm]{geometry}
\usepackage{fontspec}
\setmainfont{Latin Modern Roman}
\usepackage{amsmath,amssymb,mathtools}
\usepackage{unicode-math}
\setmathfont{Latin Modern Math}
\usepackage{xurl,hyperref}
\hypersetup{colorlinks=true,urlcolor=blue}
\setcounter{secnumdepth}{0}
\setlength{\parindent}{0pt}
\setlength{\parskip}{5pt}
\setlength{\emergencystretch}{3em}
\begin{document}
\section*{\texorpdfstring{Slice-Ribbon Conjecture}{Slice-Ribbon Conjecture}}\label{301}
\subsection{Definitions and mathematical statement}
A knot is a smooth embedding $K:S^1\hookrightarrow S^3$, considered up to ambient isotopy. It is \emph{smoothly slice} if it bounds a smooth properly embedded disk $D^2\hookrightarrow B^4$. A \emph{ribbon singularity} of an immersed disk in $S^3$ is a transverse double arc whose two preimage arcs consist of one interior arc and one properly embedded arc ending on the boundary. A ribbon knot bounds an immersed disk with only these singularities. The assertion is
\[
 K\text{ smoothly slice}\quad\Longrightarrow\quad K\text{ ribbon}.
\]
Equivalently, one asks whether every smoothly slice knot admits a slice disk with no interior local maxima for the radial Morse function, after choosing an appropriate disk. The category is smooth, not merely locally flat topological sliceness.
\par\smallskip\textit{Formulation and concept references:} \href{https://www.openproblemgarden.org/op/slice_ribbon_problem_0}{[S1]}.

\subsection{Short English statement}
Must every knot that bounds a smooth disk in four-dimensional space also bound a disk in three-dimensional space with only ribbon-type self-intersections?
\subsection{Research attempt: ribbon movies, algebraic sliceness, and maximum cancellation}
\paragraph{1. Smooth category and current-source audit.}
The question asks whether every smoothly slice knot has some ribbon disk. It does not ask that every slicing disk can be isotoped into ribbon position, and topological sliceness is not enough. Manolescu's 2026 survey [E1], Conjecture 7.2, states the smooth slice--ribbon question. Hom and Park's June 2026 paper [E2] proves a restricted ribbon classification for certain connected sums of fibered knots, not the general implication. Oliveira-Smith's March 2026 preprint [E3] gives a smoothly slice knot in the standard four-ball whose ribbon status is not known; a candidate is not a counterexample. These primary statements were checked on 27 September 2026. The attempt below supplies elementary geometric and algebraic deductions, with no claim to settle the unrestricted implication.

\paragraph{2. Ribbon implies slice: the movie makes the direction explicit.}
A ribbon presentation can be described by disjoint disk minima and embedded bands. In a three-dimensional picture, a band may pass through the interior of a disk. Each such intersection is a ribbon double arc: its preimage in the band runs between the two boundary edges, while its other preimage lies in the interior of the disk sheet. Arrange the finitely many such passages in separate neighborhoods, with no triple points or other singularities.

Place the disk minima at small distinct radial levels in $B^4$, and attach each band later in the radial coordinate. A band that appears to pass through a disk in its projection lies at a different radial level from that disk's interior. Near each ribbon arc the two sheets therefore separate in the fourth direction. After smoothing corners this gives a properly embedded surface in $B^4$, with the same boundary knot, whose radial function has only minima and saddles. If the abstract surface is a disk, it is a smooth slicing disk. This proves the easy direction and explains the absence of maxima. It does not provide a way to perform the construction from an arbitrary embedded slicing disk whose movie already contains maxima.

One can also start directly with a movie: create an unlink by minima, apply bands at later regular levels, and assume the resulting surface is connected, orientable, and has one boundary component and Euler characteristic one. The classification of compact surfaces makes it a disk. Its movie has no maximum, so it is a ribbon disk under the equivalence retained in the question.

\paragraph{3. A large constructive subcase with an exact topological check.}
Begin with $r$ disjoint disk minima. Attach $r-1$ orientable bands, each joining two previously distinct connected components, so that the incidence graph on the initial disks is a tree. The abstract surface is connected after all bands are attached. Its Euler characteristic is
\[
 \chi=r-(r-1)=1.
\]
It has one boundary component: initially every component is a disk, and joining two disks by an orientable boundary band produces another disk; induction over the tree proves the claim. Thus the resulting surface is a disk, regardless of how its bands are knotted or twisted in the ambient movie, provided the specified embeddings and orientability conditions hold. Its boundary is a ribbon knot and hence slice.

This construction gives infinitely many possible presentations, but it is not an enumeration of all slice knots. The tree records abstract connectivity and Euler characteristic; it does not record how an arbitrary knot can be obtained as the boundary of such a presentation. A proposed recognition algorithm must supply the bands and verify their boundary knot, not merely write down a tree with the correct number of edges.

The restriction that each band joins distinct components is a sufficient condition for this elementary induction. More general ribbon movies may be reorganized differently. No claim is made that a fixed diagram of every ribbon knot has a unique such tree or a bounded number of bands.

\paragraph{4. Euler characteristic alone does not cancel maxima.}
Put a smooth slicing disk in generic radial Morse position. Let $c_0,c_1,c_2$ be the numbers of interior minima, saddles, and maxima. Since its Euler characteristic is one,
\[
 c_0-c_1+c_2=1.
\]
For a ribbon disk, $c_2=0$ and $c_1=c_0-1$. For a general disk, the equality only says $c_1=c_0+c_2-1$. It does not pair each maximum with a geometrically cancellable saddle.

Morse cancellation requires an appropriate connecting trajectory and control of the attaching data. An algebraic cancellation in a chain complex is weaker than a geometric cancellation of handles. In four-dimensional disk complements, knotting and framing information remain after homology has been simplified. Adding a cancelling critical-point pair may make a diagram easier to manipulate, but it also changes the diagram and does not imply that all maxima can eventually be removed.

The quantifier ``some disk'' matters here. Finding a slicing disk whose given handle decomposition resists an attempted cancellation does not prove that its boundary knot is not ribbon. Another disk for the same knot may have a completely different complement or Morse movie. A counterexample needs an obstruction to every ribbon presentation of that knot.

\paragraph{5. The Fox--Milnor algebraic condition derived from a metabolizer.}
Use the classical theorem that a smoothly slice knot is algebraically slice: some Seifert form has a rank-$g$ direct summand on which the form vanishes, where the Seifert surface has genus $g$. This theorem is an external knot-concordance input. Once such a metabolizer is given, the polynomial consequence is elementary.

Choose a basis beginning with the metabolizer. The $2g\times2g$ Seifert matrix has block form
\[
 V=\begin{pmatrix}0&A\\ B&C\end{pmatrix}.
\]
The Alexander polynomial is represented, up to a unit $\pm t^m$, by $\det(V-tV^T)$. Expanding the determinant of the zero-upper-left block matrix gives
\[
 \det(V-tV^T)=(-1)^g\det(A-tB^T)\det(B-tA^T).
\]
Let $f(t)=\det(A-tB^T)$. Transposition and factoring $-t$ from each row show
\[
 \det(B-tA^T)=(-t)^g f(t^{-1}).
\]
Hence
\[
 \Delta_K(t)=\pm t^m f(t)f(t^{-1})
\]
for an integer exponent $m$, with normalization absorbed into the allowed unit. Evaluating at $t=-1$ gives $|\Delta_K(-1)|=|f(-1)|^2$, a square integer. Therefore a knot with nonsquare determinant cannot be smoothly slice and cannot be ribbon.

The implication runs from a metabolizer to a factorization. The determinant factorization alone does not produce an embedded disk, and certainly does not produce a ribbon movie. Both slice and ribbon knots satisfy this same necessary condition, so it cannot distinguish the two classes by itself.

\paragraph{6. An exact genus-one stress test.}
Consider a Seifert matrix of the form
\[
 V=\begin{pmatrix}0&a\\a-1&b\end{pmatrix},\qquad a,b\in\mathbb Z.
\]
Its skew-symmetrization is $\left(\begin{smallmatrix}0&1\\-1&0\end{smallmatrix}\right)$, as required for a genus-one Seifert form, and the first basis vector spans a metabolizer. Direct calculation yields
\[
 \det(V-tV^T)
 =-\bigl(a-(a-1)t\bigr)\bigl((a-1)-at\bigr).
\]
The result is independent of $b$, while the matrix itself varies with $b$. Its determinant at $t=-1$ has absolute value $(2a-1)^2$. For $a=2$, the polynomial is $-2+5t-2t^2$, equivalent to the symmetric normalization $5-2(t+t^{-1})$, and the knot determinant is nine.

This calculation constructs and tests algebraic forms. It does not identify every realization of those forms as a slice or ribbon knot. The invisibility of $b$ to this Alexander polynomial is a concrete warning that such a polynomial records only part of the geometric information. Claiming that the square determinant or displayed factorization proves sliceness would reverse the established implication.

The ordinary knot signature also vanishes for a nonsingular symmetric form with a half-dimensional isotropic subspace. Indeed such a subspace forces the positive and negative indices each to be at least $g$, and their sum is $2g$, so both equal $g$. Applied to the appropriate symmetrized Seifert form, this supplies another shared slice obstruction, not a ribbon-versus-slice separation.

\paragraph{7. Disk complements explain another tempting false shortcut.}
A ribbon movie with $r$ minima and $r-1$ bands gives a handle description of its complement with one-handles corresponding to minima and two-handles corresponding to bands. The associated meridian presentation has $r$ generators and $r-1$ band relations. Each band relation identifies one meridian with a conjugate of another; in the abelianization it identifies the corresponding generators. For the connected tree presentation of Section 3, all meridians become equal, giving abelianization $\mathbb Z$.

A general slicing disk also has complement with first homology $\mathbb Z$, by the usual meridian/Alexander-duality calculation. Its maxima contribute higher-index handles. Those handles do not change the fundamental group presentation in the same way as one- and two-handles, yet they may be geometrically essential in that particular decomposition. Thus agreement of $H_1$, or even a promising presentation of $\pi_1$, does not supply a cancellation of the higher handles. Group-theoretic simplifications need a justified geometric realization with the needed framings.

Likewise, a locally flat topological disk obtained by a topological surgery theorem is not automatically a smooth disk in the standard $B^4$. A proof that crosses categories at this point no longer addresses the supplied implication. The recent candidate in [E3] is relevant precisely because its slicing disk is in the standard smooth four-ball, but its being slice does not itself establish nonribbonness.

\paragraph{8. Diagnostics and outcome.}
The saved exact checks enumerate finite tree incidence data, verify the Euler counts, and compute the genus-one Alexander polynomials and square determinants over the integers. They do not recognize arbitrary knot isotopy classes, search all band presentations, or certify a nonribbon knot. The ribbon construction and the algebraic necessary conditions are justified by the arguments above and the explicitly named standard input.

\textbf{UNRESOLVED.} The attempt proves a concrete ribbon construction and derives exact shared algebraic obstructions from a metabolizer. The first unrestricted gap is a geometric procedure replacing an arbitrary smooth slicing disk by some disk with no radial maxima, or an invariant obstructing every such disk for a particular slice knot. Critical-point counts and Alexander-polynomial factorization do not fill that gap.

\subsection{Sources}
\begin{itemize}
\item[S1]Fox, R. H. "Some problems in knot theory," in Topology of 3-manifolds and related topics (Proc. The Univ. of Georgia Institute, 1961), pp. 168-176, 1962.
\url{https://www.openproblemgarden.org/op/slice_ribbon_problem_0}.
\textit{Formulation source.}
\item[E1]C. Manolescu, From knots to four-manifolds, arXiv:2601.05425, Section 7 and Conjecture 7.2.
\url{https://arxiv.org/html/2601.05425}.
\textit{Primary 2026 survey checked 27 September 2026 for the smooth ribbon implication and its converse question.}
\item[E2]J. Hom and J. Park, Ribbon knots and iterated cables of fibered knots, Mathematische Zeitschrift 313 (2026), article 53.
\url{https://link.springer.com/article/10.1007/s00209-026-04050-3}.
\textit{Published 18 June 2026; primary abstract and introduction checked 27 September 2026. Restricted fibered-knot result, not the full conjecture.}
\item[E3]T. Oliveira-Smith, A Dunfield--Gong 4-Sphere is Standard, arXiv:2603.23717.
\url{https://arxiv.org/abs/2603.23717}.
\textit{Primary abstract checked 27 September 2026. The stated slice knot has unknown ribbon status and is only a potential counterexample.}
\end{itemize}
\end{document}
